% ============================================================
% Lettre de motivation — gabarit
% Expéditeur en haut à gauche, destinataire en haut à droite,
% date, objet, puis trois paragraphes : Vous, Je, Nous.
% A4 portrait, Roboto 11 pt, une page. Compile avec tectonic.
%
% Tout ce qui est entre <chevrons> se remplace ; rien d'autre
% ne change. Dans le texte : € s'écrit \texteuro{}, % s'écrit
% \%, & s'écrit \&.
% La ligne du destinataire part avec lui quand l'annonce ne
% nomme personne.
% ============================================================
\documentclass[11pt,a4paper]{article}
\usepackage[utf8]{inputenc}
\usepackage[T1]{fontenc}
\usepackage[sfdefault,type1]{roboto}
\renewcommand{\familydefault}{\sfdefault}
\usepackage[french]{babel}
\usepackage[a4paper,left=22mm,right=22mm,top=20mm,bottom=20mm]{geometry}
\usepackage{xcolor}
\usepackage{microtype}
\usepackage{hyperref}

\definecolor{ink}{HTML}{1A1D21}
\definecolor{body}{HTML}{3C4247}

\pagestyle{empty}
\hypersetup{hidelinks}
\setlength{\parindent}{0pt}
\setlength{\parskip}{10pt}
\linespread{1.12}
\color{body}

\begin{document}

% ---- Expéditeur (gauche) et destinataire (droite) ----
\begin{minipage}[t]{0.5\textwidth}
  {\Large\color{ink}\robotobold <Prénom Nom>}\\[4pt]
  <Rue>\\
  <Code postal Ville>\\[4pt]
  <Mobile>\\
  \href{mailto:<Email>}{<Email>}
\end{minipage}%
\hfill
\begin{minipage}[t]{0.4\textwidth}
  {\color{ink}\robotobold <Société>}\\
  <Destinataire>\\
  <Adresse de la société>\\
  <Code postal Ville de la société>
\end{minipage}

\vspace{18pt}
\begin{flushright}
<Ville>, le <date>
\end{flushright}

\vspace{6pt}
{\color{ink}\robotobold Objet : Lettre de motivation pour le poste de <Poste>}

\vspace{4pt}
Madame, Monsieur,

<Paragraphe Vous>

<Paragraphe Je>

<Paragraphe Nous>

Cordialement,

\vspace{14pt}
{\robotomedium <Prénom Nom>}

\end{document}
